% Chapter 5 Results and Evaluation   (we have some choices here)

% Illustration the system – few sample run throughs with sample queries and recommendations -tying back to the design and the various components. 
% This will show the reader how the system works.

% Evaluation

% Talk about the classic metrics ( can cover in Chap 2) and they limitations in complex system. 
% User -evaluation – tasks for users to undertake. Questions asked – we need to design these; Objective measure of the recommendations; Subjective feedback.

% Efficiency issues


Evaluation of the project will be done in four parts:
\begin{itemize}
    \item Recommendation System Evaluation
    \item Large Language Model Evaluation
    \item Feedback and Recommendation Adjustment Evaluation
    \item User Evaluation
\end{itemize}

\section{Recommendation System Evaluation}
\label{sec:recommendation_system_evaluation}
As mentioned in \ref{sec:recommendation_system}, the recommendation system is designed to be a proof of concept, and is not intended to be the best implementation.

What the system is being evaluated on is:
\begin{itemize}
    \item The system works reliably
    \item That the system is able to provide recommendations that make sense
    \item The recommendation breakdown is generated and is understandable
    \item The system is able to provide recommendations based on the features
    \item The system provides recommendations unique to user preferences
\end{itemize}

As mentioned in 
As this the recommendation system is only there to provide a working implementation of a recommendation system, with adjustable parameters, it is not the main focus of the project.
As such, it is not intended to be necessarily the best recommendation system, there is little need to evaluate the recommendation system in detail.


\section{Large Language Model Evaluation}
\label{sec:large_language_model_evaluation}
The large language model is evaluated on the following:
\begin{itemize}
    \item The agent can work with the recommendation system
    \item The agent works with the user to explore their preferences
    \item The agent can provide a breakdown of the recommendations with an explanation that is understandable
    \item The agent can handle out of scope conversations and can redirect the conversation back to the task at hand
    \item The system is timely and responsive
\end{itemize}


\section{User Evaluation}
\label{sec:user_evaluation}
For this project, as classic metrics do not apply, we need to come up with a new way to evaluate the system. 
We need to design a user evaluation, where users undertake tasks and provide feedback. This will provide an objective measure of the recommendations and subjective feedback.


The user questions will be designed to evaluate the system as a whole, as well as the recommendation system and the large language model.
The user will be asked to undertake one of two scenarios, depending on their preferences.

\subsection{User Scenarios}
\label{sec:user_scenarios}
The users will be asked to undertake one of two scenarios, depending on their preferences.
\paragraph{Scenario 1:}
\begin{enumerate}
    \item The user will be asked to start with a blank profile and talk to the agent about their film preferences on topics such as films, genres, directors, etc.
    \item The system will provide a set of recommendations based on the user preferences.
    \item The user will evaluate the given recommendations and provide feedback on the recommendations.
    \item The user will ask the system to provide a breakdown and explanation of the recommendations.
    \item The user will discuss with the agent what feature are important to them in recommendations and generate a new set of recommendations based on those alterations.
    \item Finally, the user will provide feedback on the system, the recommendations, and the system as a whole.
\end{enumerate}

\paragraph{Scenario 2:}
\begin{enumerate}
    \item The user will be asked to start with a template sample profile.
    \item The system will provide a set of recommendations based on the user preferences.
    \item The user will ask the system to provide a breakdown and explanation of the recommendations.
    \item The user will provide feedback on the recommendations and the explanation.
    \item The user will discuss with the agent what feature are important to them in recommendations and generate a new set of recommendations based on those alterations.
    \item The user will evaluate the given recommendations and provide feedback on the recommendations.
    \item Finally, the user will provide feedback on the system, the recommendations, and the system as a whole.
\end{enumerate}


\subsection{User Survey}
\label{sec:user_survey}
Using the following questions, users will be asked to provide feedback on the system, the recommendations, and the system as a whole
\begin{itemize}
    \item Does the system feel useful
    \item How cumbersome is it to use the system
    \item How well does the system understand your preferences
    \item How useful is the recommendations explanation
    \item Did the system reflect your change in preferences in the recommendations
    \item Would you use the system again
    \item What do you like about the system
    \item What do you dislike about the system
    \item What would you like to see added to the system
    \item What do you think of the system as a whole
    \item How would you rate the systems objective overall
\end{itemize}



